\scanpage{18}

2.3. Considérons \unsure{à nouveau|: à nouveau} l'iso. can.

\[
\underline{\mathrm{End}}(E)/\underline{\mathcal O}\cdot1_E\simeq\underline{\mathrm{Quad}}(E)\otimes\underline L\;(\simeq\mathrm{Quad}(E,L))
\]

Soit $u$ une section de $\underline{\mathrm{End}}(E)$ i.e. un endom. de $E$, \deleted{considér\unsure{ons}} son image est donc une forme quadratique $Q_u$ à valeurs dans $L$. Je dis qu'on a

\[
\boxed{Q_u(x)=u(x)\wedge x}
\]

(\deleted{mais} il faut vérifier le signe !). \deleted{C'est} On revérifie directement par des coordonnées que

\[
Q_u=0\quad\Longleftrightarrow\quad u\in\Gamma\underline{\mathcal O}\cdot\mathrm{id}_E
\]

[ NB. Si

\[
u=\begin{pmatrix}a&b\\c&d\end{pmatrix}
\]

alors

\[
Q_u(\lambda,\mu)=c\lambda^2+(d-a)\lambda\mu\deleted{\unsure{+|-}}\unsure{-|+}b\mu^2
\]

\deleted{qui} est id. nul en $\lambda,\mu$ ssi $b=c=0$, $d=a$ i.e. $u=\deleted{\unsure{a}}\begin{pmatrix}a&\deleted{\unsure{c|a}}0\\0&a\end{pmatrix}\deleted{\unsure{=a}}=a\,\mathrm{id}_E$

Si $Q=\alpha\lambda^{\unsure{2|7}}+\beta\lambda\mu+\gamma\mu^2$, alors les $u$ tels que $Q_u=Q$ sont ceux de la forme

\[
u_a=\begin{pmatrix}a&-\gamma\\\alpha&a+\beta\end{pmatrix},\quad\text{avec }a\in\Gamma\underline{\mathcal O}\text{ arbitraire}
\]

parmi lesquels

\[
u_0=\begin{pmatrix}0&-\gamma\\\alpha&\beta\end{pmatrix}\qquad\ldots\qquad]
\]

On a \deleted{alors}
% archive sheet number: 17
